% !TEX root =./ECE444_Practice_main.tex
%
% L13 practice set (Measurement Lab 1 - Impedance and S-parameters). Every
% problem is solvable from given data: no VNA required.

\fullwidth{\textbf{Learning Objective 2.6: } I can measure the impedance and S-parameters
of an antenna using a vector network analyzer and interpret the results.
\\[4pt]\normalfont\small Show your work. A \textbf{Documentation} line at the end
is required (write \textbf{None} if you did not collaborate).}

\begin{questions}

%%%%%%%%%%%%%%%%%%%%%%%%%%%%%%%%%%%%%%%%%%%%%%%%%%%%%%%%%%%%%%%%%%%%%%%%%%%%
\question \textbf{Four names for one number.} A one-port measurement gets reported
four different ways depending on which menu you open. Reference impedance is
$Z_0 = 50$~\ohms throughout.

\begin{parts}

	\begin{minipage}[t]{0.58\linewidth}
		\part An antenna reads $|S_{11}| = -14\db$ at its design frequency.
		Find $|\Gamma|$, the VSWR, and the percentage of incident power reflected.
		\begin{solution}[1.2in]
		\eq{
			|\Gamma| &= 10^{-14/20} = 0.199 \\
			\text{VSWR} &= \frac{1+0.199}{1-0.199} = \frac{1.199}{0.801} = 1.50 \\
			\frac{P_{refl}}{P_{inc}} &= |\Gamma|^2 = 0.0398 \rightarrow 4.0\%
		}
		\end{solution}
	\end{minipage}\hfill%
	\begin{minipage}[t]{0.37\linewidth}
		\raggedleft \vspace{12pt}
		$|\Gamma|$ = \ansbox[1.1in]{$0.199$} \\[8pt]
		VSWR = \ansbox[1.1in]{$1.50$} \\[8pt]
		$P_{refl}$ = \ansbox[1.1in]{$4.0\%$}
	\end{minipage}

	\begin{minipage}[t]{0.58\linewidth}
		\part A second antenna shows VSWR $= 2.0$ at the same frequency.
		What is its $|S_{11}|$ in dB?
		\begin{solution}[1.0in]
		\eq{
			|\Gamma| &= \frac{\text{VSWR}-1}{\text{VSWR}+1} = \frac{1}{3} = 0.333 \\
			|S_{11}| &= 20\ \mylog{0.333} = -9.5\db
		}
		It misses the $-10\db$ specification, but only just.
		\end{solution}
	\end{minipage}\hfill%
	\begin{minipage}[t]{0.37\linewidth}
		\raggedleft \vspace{12pt}
		$|S_{11}|$ = \ansbox[1.3in]{$-9.5\db$}
	\end{minipage}

	\part Why is $-10\db$ the conventional acceptance level for an antenna match?
	Answer in terms of power, in one or two sentences.
		\begin{solution}[0.9in]
		At $-10\db$, $|\Gamma|^2 = 0.10$: exactly $10\%$ of the delivered power comes
		back and $90\%$ goes into the antenna. Below that level the reflected power is
		a small correction to the link budget, so making the match better buys very
		little, while making it worse costs quickly.
		\end{solution}

	\begin{minipage}[t]{0.58\linewidth}
		\part Lesson 7 predicted $Z_{\text{in}} = 73 + j42.5$~\ohms for an ideal half-wave
		dipole. Compute $|\Gamma|$, $|S_{11}|$ in dB, and the VSWR. Does the textbook
		dipole meet the $-10\db$ specification, and what one change fixes it?
		\begin{solution}[1.5in]
		\eq{
			\Gamma &= \frac{23 + j42.5}{123 + j42.5} \\
			|\Gamma| &= \frac{48.3}{130.1} = 0.371 \\
			|S_{11}| &= 20\ \mylog{0.371} = -8.6\db \\
			\text{VSWR} &= \frac{1.371}{0.629} = 2.18
		}
		It \textbf{fails}. The fix is to shorten the element to about $0.47$--$0.48\lambda$
		so the reactance goes to zero and $R_{\text{in}}$ falls to roughly $70$~\ohms:
		\eq{
			|\Gamma| = \frac{20}{120} = 0.167 \rightarrow -15.6\db, \qquad \text{VSWR} = 1.40
		}
		The $42.5$~\ohms of reactance is the entire difference between a marginal
		antenna and a good one.
		\end{solution}
	\end{minipage}\hfill%
	\begin{minipage}[t]{0.37\linewidth}
		\raggedleft \vspace{12pt}
		$|S_{11}|$ = \ansbox[1.3in]{$-8.6\db$} \\[8pt]
		VSWR = \ansbox[1.3in]{$2.18$} \\[8pt]
		Passes? \ansbox[1.3in]{no}
	\end{minipage}

\end{parts}

\newpage
%%%%%%%%%%%%%%%%%%%%%%%%%%%%%%%%%%%%%%%%%%%%%%%%%%%%%%%%%%%%%%%%%%%%%%%%%%%%
\question \textbf{From a marker readout to an impedance.} A calibrated VNA marker
sitting at $1.575\ \ghz$ on a small GPS-band element reads
$S_{11} = 0.33\ \angle 155\mydeg$.

\begin{parts}

	\begin{minipage}[t]{0.58\linewidth}
		\part Write $\Gamma$ in rectangular form and find the input impedance $Z_{\text{in}}$.
		\begin{solution}[1.6in]
		\eq{
			\Gamma &= 0.33\lp \cos 155\mydeg + j \sin 155\mydeg \rp = -0.299 + j0.140 \\
			Z_{\text{in}} &= Z_0\ \frac{1+\Gamma}{1-\Gamma}
			     = 50\ \frac{0.701 + j0.140}{1.299 - j0.140}
		}
		Multiplying numerator and denominator by the conjugate of the denominator,
		\eq{
			Z_{\text{in}} &= 50\ \frac{0.891 + j0.279}{1.707} \\
			       &= 50 \lp 0.522 + j0.163 \rp \\
			       &= 26.1 + j8.2\ \ohms
		}
		\end{solution}
	\end{minipage}\hfill%
	\begin{minipage}[t]{0.37\linewidth}
		\raggedleft \vspace{12pt}
		$Z_{\text{in}}$ = \ansbox[1.5in]{$26.1 + j8.2$~\ohms}
	\end{minipage}

	\part The reactance came out positive. Is the element electrically long or short
	at $1.575\ \ghz$, is its resonance above or below that frequency, and which way
	would you trim it?
		\begin{solution}[1.0in]
		Positive reactance is inductive, which is what a wire antenna looks like when
		it is \textbf{electrically long}. Its resonance ($X = 0$) therefore lies
		\textbf{below} $1.575\ \ghz$. To move the resonance up onto the design
		frequency you trim the element \textbf{shorter}.
		\end{solution}

	\begin{minipage}[t]{0.58\linewidth}
		\part Find $|S_{11}|$ in dB and the VSWR. Does this element meet the
		$-10\db$ specification as measured?
		\begin{solution}[1.1in]
		\eq{
			|S_{11}| &= 20\ \mylog{0.33} = -9.6\db \\
			\text{VSWR} &= \frac{1.33}{0.67} = 1.99
		}
		No -- it misses the specification by $0.4\db$.
		\end{solution}
	\end{minipage}\hfill%
	\begin{minipage}[t]{0.37\linewidth}
		\raggedleft \vspace{12pt}
		$|S_{11}|$ = \ansbox[1.3in]{$-9.6\db$} \\[8pt]
		VSWR = \ansbox[1.3in]{$1.99$}
	\end{minipage}

	\part Suppose you trim the element perfectly, so at $1.575\ \ghz$ the reactance
	is zero and the resistance stays near $26$~\ohms. Will the match pass then?
	Compute it, and state what the result tells you about relying on trimming alone.
		\begin{solution}[1.4in]
		\eq{
			\Gamma = \frac{26-50}{26+50} = -0.316 \quad\rightarrow\quad
			|S_{11}| = -10.0\db, \qquad \text{VSWR} = 1.92
		}
		It lands \emph{exactly} on the specification with no margin at all. Trimming
		only removes reactance; it cannot move the resistance. An element whose
		radiation resistance sits far from $50$~\ohms needs a matching network or a
		different feed arrangement, not a shorter piece of wire.
		\end{solution}

\end{parts}

\newpage
%%%%%%%%%%%%%%%%%%%%%%%%%%%%%%%%%%%%%%%%%%%%%%%%%%%%%%%%%%%%%%%%%%%%%%%%%%%%
\question \textbf{Bandwidth off a sweep.} A calibrated, load-verified sweep of a
wire dipole produced the following magnitudes.

\begin{center}
\begin{tabular}{l|cccccccccc}
$f$ (MHz)      & 850 & 870 & 880 & 890 & 900 & 905 & 910 & 920 & 930 & 950 \\ \hline
$|S_{11}|$ (dB) & $-3.8$ & $-6.5$ & $-8.8$ & $-12.4$ & $-17.7$ & $-19.4$ & $-17.8$ & $-12.5$ & $-9.0$ & $-5.2$
\end{tabular}
\end{center}

\begin{parts}

	\begin{minipage}[t]{0.58\linewidth}
		\part Identify the resonant frequency and say how you know.
		\begin{solution}[0.9in]
		The deepest point of the sweep is $-19.4\db$ at $905\ \mhz$, and the trace is
		symmetric about it. That is the resonance: the reactance passes through zero
		there, so the reflection is smallest.
		\end{solution}
	\end{minipage}\hfill%
	\begin{minipage}[t]{0.37\linewidth}
		\raggedleft \vspace{12pt}
		$f_0$ = \ansbox[1.2in]{$905\ \mhz$}
	\end{minipage}

	\begin{minipage}[t]{0.58\linewidth}
		\part Interpolate the two $-10\db$ crossings and report the impedance
		bandwidth both in MHz and as a percentage.
		\begin{solution}[1.5in]
		Lower crossing, between $880\ \mhz$ ($-8.8\db$) and $890\ \mhz$ ($-12.4\db$):
		\eq{
			f_1 = 880 + 10\ \frac{10-8.8}{12.4-8.8} = 883\ \mhz
		}
		Upper crossing, between $920\ \mhz$ ($-12.5\db$) and $930\ \mhz$ ($-9.0\db$):
		\eq{
			f_2 = 920 + 10\ \frac{12.5-10}{12.5-9.0} = 927\ \mhz
		}
		\eq{
			\text{BW} = 927 - 883 = 44\ \mhz, \qquad
			\frac{44}{905} = 4.8\%
		}
		\end{solution}
	\end{minipage}\hfill%
	\begin{minipage}[t]{0.37\linewidth}
		\raggedleft \vspace{12pt}
		BW = \ansbox[1.2in]{$44\ \mhz$} \\[8pt]
		Fractional = \ansbox[1.2in]{$4.8\%$}
	\end{minipage}

	\begin{minipage}[t]{0.58\linewidth}
		\part At resonance the Smith-chart marker sits on the real axis, to the
		\emph{right} of centre. Find $Z_{\text{in}}$ at resonance.
		\begin{solution}[1.2in]
		On the real axis $\Gamma$ is purely real, and right of centre means it is
		positive, so $R_{\text{in}} > 50$~\ohms:
		\eq{
			|\Gamma| &= 10^{-19.4/20} = 0.107 \\
			Z_{\text{in}} &= 50\ \frac{1+0.107}{1-0.107} = 50 \lp 1.240 \rp = 62\ \ohms
		}
		Without the Smith chart you could not choose between $62$~\ohms and
		$40$~\ohms: the magnitude alone does not carry the sign.
		\end{solution}
	\end{minipage}\hfill%
	\begin{minipage}[t]{0.37\linewidth}
		\raggedleft \vspace{12pt}
		$Z_{\text{in}}$ = \ansbox[1.4in]{$62 + j0$~\ohms}
	\end{minipage}

	\part A thin wire dipole typically shows $3$--$10\%$ impedance bandwidth. Does
	this element fit the rule of thumb, and what physical change would widen it?
		\begin{solution}[1.0in]
		Yes: $4.8\%$ sits in the lower half of the usual $3$--$10\%$ range, which is
		what a thin conductor gives. Bandwidth widens when the antenna's $Q$ drops, so
		use a fatter conductor -- a thicker rod, a sleeve, a cage of wires, or a
		biconical element -- which lowers the stored-to-radiated energy ratio.
		\end{solution}

\end{parts}

\newpage
%%%%%%%%%%%%%%%%%%%%%%%%%%%%%%%%%%%%%%%%%%%%%%%%%%%%%%%%%%%%%%%%%%%%%%%%%%%%
\question \textbf{Calibration and the reference plane.} A classmate calibrates
short--open--load at the end of the test cable, then screws a $10\ \text{cm}$ RG-58
pigtail ($v_f = 0.66$) between that connector and the antenna before sweeping.

\begin{parts}

	\part Name the three error terms a one-port SOL calibration removes, one
	sentence each, and explain why exactly three standards are required.
		\begin{solution}[1.5in]
		\textbf{Directivity}: the coupler that should sample only the returning wave
		leaks some of the outgoing wave into the same receiver, so the VNA sees a
		reflection even from a perfect load. \textbf{Source match}: the test port is
		not exactly $50$~\ohms, so energy returned by the antenna is partly
		re-reflected back at it. \textbf{Tracking}: the reference and test receiver
		paths have different gain and phase, and both drift with frequency. Three
		unknowns require three independent known loads, which is exactly what short
		($\Gamma = -1$), open ($\Gamma = +1$), and load ($\Gamma = 0$) provide.
		\end{solution}

	\part Describe what the extra pigtail does to the measured $|S_{11}|$ magnitude
	and, separately, to the measured impedance. Why is this combination dangerous?
		\begin{solution}[1.3in]
		A short low-loss pigtail changes the \emph{magnitude} hardly at all, so the
		dB plot still looks correct and the resonant frequency still reads correctly.
		It adds round-trip \emph{phase}, which rotates the entire trace around the
		Smith chart, so the impedance readout is wrong -- possibly wildly wrong,
		including the sign of the reactance. That combination is dangerous because
		the plot that still looks correct is the one most likely to be trusted.
		\end{solution}

	\begin{minipage}[t]{0.58\linewidth}
		\part At $915\ \mhz$, find the guided wavelength in the pigtail, its length
		in wavelengths, and the round-trip phase shift $2\beta\ell$ it adds.
		\begin{solution}[1.4in]
		\eq{
			\lambda_g &= \frac{c\ v_f}{f} = \frac{\lp 3\times10^8 \rp \lp 0.66 \rp}{915\times10^6}
			 = 0.216\ \m = 21.6\ \text{cm} \\
			\frac{\ell}{\lambda_g} &= \frac{10}{21.6} = 0.462 \\
			2\beta\ell &= 2 \lp 360\mydeg \rp \lp 0.462 \rp = 333\mydeg
		}
		That is nearly a full rotation of the chart.
		\end{solution}
	\end{minipage}\hfill%
	\begin{minipage}[t]{0.37\linewidth}
		\raggedleft \vspace{12pt}
		$\lambda_g$ = \ansbox[1.2in]{$21.6\ \text{cm}$} \\[8pt]
		$2\beta\ell$ = \ansbox[1.2in]{$333\mydeg$}
	\end{minipage}

	\begin{minipage}[t]{0.58\linewidth}
		\part At what frequency is this pigtail exactly a half guided wavelength
		long? Why does that frequency matter to your classmate?
		\begin{solution}[1.3in]
		\eq{
			\lambda_g &= 2\ell = 0.20\ \m \\
			f &= \frac{c\ v_f}{\lambda_g} = \frac{\lp 3\times10^8 \rp \lp 0.66 \rp}{0.20} \\
			  &= 990\ \mhz
		}
		A half-wavelength line repeats its load impedance, so at $990\ \mhz$ (and
		multiples of it) the un-de-embedded reading happens to be correct. At every
		other frequency it is rotated. Obtaining the correct answer at that one
		frequency is a coincidence rather than a calibration; fix it with port extension
		or by calibrating at the antenna connector.
		\end{solution}
	\end{minipage}\hfill%
	\begin{minipage}[t]{0.37\linewidth}
		\raggedleft \vspace{12pt}
		$f$ = \ansbox[1.2in]{$990\ \mhz$}
	\end{minipage}

\end{parts}

\newpage
%%%%%%%%%%%%%%%%%%%%%%%%%%%%%%%%%%%%%%%%%%%%%%%%%%%%%%%%%%%%%%%%%%%%%%%%%%%%
\question \textbf{What the VNA cannot see.} The instrument reports one complex
number per frequency. This problem is about the questions that number cannot answer.

\begin{parts}

	\part You solder a $50$~\ohms chip resistor across the connector and sweep it.
	Predict $|S_{11}|$, the VSWR, and the radiated power. What does the result prove
	about what $S_{11}$ measures?
		\begin{solution}[1.2in]
		$\Gamma = 0$, so $|S_{11}|$ drops to the noise floor of the instrument and the
		VSWR reads $1.00$ at every frequency: a perfect match. It radiates
		\textbf{nothing} -- all of the delivered power becomes heat. $S_{11}$ therefore
		measures \textbf{mismatch only}. It cannot separate power that left as
		radiation from power that died as loss, so it can never confirm efficiency.
		\end{solution}

	\begin{minipage}[t]{0.58\linewidth}
		\part Antenna A measures $-25\db$ with $2.0\dbi$ of gain; antenna B measures
		$-9\db$ with $5.0\dbi$. Find the mismatch loss of each and decide which one
		puts more power into its main beam.
		\begin{solution}[1.6in]
		\eq{
			|\Gamma_A| &= 10^{-25/20} = 0.056, \qquad |\Gamma_A|^2 = 0.0032 \\
			L_A &= -10\ \mylog{1 - 0.0032} = 0.01\db \\
			|\Gamma_B| &= 10^{-9/20} = 0.355, \qquad |\Gamma_B|^2 = 0.126 \\
			L_B &= -10\ \mylog{1 - 0.126} = 0.58\db
		}
		Realized gains are $2.0 - 0.01 = 2.0\dbi$ and $5.0 - 0.58 = 4.4\dbi$.
		\textbf{Antenna B wins by about $2.4\db$}, despite its much worse $S_{11}$.
		Mismatch loss is small until $|S_{11}|$ gets much worse than $-10\db$.
		\end{solution}
	\end{minipage}\hfill%
	\begin{minipage}[t]{0.37\linewidth}
		\raggedleft \vspace{12pt}
		$L_A$ = \ansbox[1.2in]{$0.01\db$} \\[8pt]
		$L_B$ = \ansbox[1.2in]{$0.58\db$} \\[8pt]
		Better: \ansbox[1.2in]{B}
	\end{minipage}

	\part During the lab you bring a hand within a few centimetres of the element
	and the dip moves. State which way the resonant frequency shifts, what happens
	to the depth of the dip, and why neither effect is instrument error.
		\begin{solution}[1.2in]
		The resonance shifts \textbf{down} and the dip usually gets \textbf{shallower
		and broader}. Tissue is a lossy dielectric: placed in the reactive near field
		it adds capacitive loading, which lowers $f_0$, and it adds loss, which lowers
		the antenna's $Q$. The input impedance changed, so the VNA is
		reporting a correct measurement -- the hand became part of the antenna.
		\end{solution}

	\part Your antenna measures $-30\db$ at the design frequency, yet the link will
	not close. Name two things a VNA cannot see that would explain this, and say how
	you would check one of them.
		\begin{solution}[1.3in]
		(1) \textbf{Efficiency}: a lossy structure -- corroded joints, a resistive
		loading element, wet material in the near field -- absorbs the power it
		accepts while still presenting a good match. (2) \textbf{Pattern or
		polarization}: the beam may point somewhere other than the far end of the
		link, or be cross-polarized to it. Check the first with a gain comparison
		against a standard-gain antenna on the range (Lesson 14): if the measured gain
		falls well short of the directivity you predicted, the missing decibels are
		loss.
		\end{solution}

\end{parts}

\vspace{1.5em}
\noindent\textbf{Documentation:} \rule[-2pt]{0.6\linewidth}{0.4pt}

\end{questions}
